\chapter*{Webographie}
\addcontentsline{toc}{chapter}{Webographie}
\markboth{Webographie}{Webographie}

\vspace*{0.5cm}

\begin{enumerate}

\item \textbf{Groq — Documentation officielle de l'API}\\
Documentation complète de l'API Groq et du modèle Llama 3.3-70b.\\
\url{https://console.groq.com/docs}\\
\textit{Consulté en mai 2026}

\item \textbf{React.js — Documentation officielle}\\
Guide complet de React.js pour la création d'interfaces utilisateur.\\
\url{https://react.dev}\\
\textit{Consulté en mai 2026}

\item \textbf{Flask — Documentation officielle}\\
Documentation du micro-framework web Python Flask.\\
\url{https://flask.palletsprojects.com}\\
\textit{Consulté en mai 2026}

\item \textbf{Monaco Editor — Documentation}\\
Documentation de l'éditeur de code Monaco (moteur de VS Code).\\
\url{https://microsoft.github.io/monaco-editor}\\
\textit{Consulté en mai 2026}

\item \textbf{StarUML — Documentation officielle}\\
Guide d'utilisation de StarUML et export XMI.\\
\url{https://docs.staruml.io}\\
\textit{Consulté en mai 2026}

\item \textbf{Vercel — Documentation de déploiement}\\
Guide de déploiement d'applications React sur Vercel.\\
\url{https://vercel.com/docs}\\
\textit{Consulté en mai 2026}

\item \textbf{Render — Documentation de déploiement}\\
Guide de déploiement d'applications Python Flask sur Render.\\
\url{https://render.com/docs}\\
\textit{Consulté en mai 2026}

\item \textbf{PyMuPDF — Documentation officielle}\\
Documentation de la bibliothèque PyMuPDF pour le traitement de PDF.\\
\url{https://pymupdf.readthedocs.io}\\
\textit{Consulté en mai 2026}

\item \textbf{UML — Spécification OMG}\\
Spécification officielle du langage UML par l'Object Management Group.\\
\url{https://www.omg.org/spec/UML}\\
\textit{Consulté en mai 2026}

\item \textbf{XMI — Spécification OMG}\\
Spécification officielle du format XMI par l'Object Management Group.\\
\url{https://www.omg.org/spec/XMI}\\
\textit{Consulté en mai 2026}

\item \textbf{GitHub — Dépôt du projet UML-to-Code}\\
Code source complet du projet disponible en open source.\\
\url{https://github.com/ABAKAR5/uml-to-code}\\
\textit{Consulté en mai 2026}

\item \textbf{Application déployée — UML-to-Code}\\
Application web déployée et accessible en ligne.\\
\url{https://uml-to-code-peach.vercel.app}\\
\textit{Consulté en mai 2026}

\end{enumerate}